\section{Introduction}

Financial QA combines requirements that are individually difficult and jointly unforgiving: locating the relevant period and entity, interpreting tables and prose, executing arithmetic, preserving units and signs, and grounding narrative conclusions in filings. Iterative language-model agents are a natural response. A ReAct-style system can search, inspect evidence, revise a hypothesis, and stop when it has an answer. That flexibility, however, has a direct operational cost. Every additional step consumes tokens and latency, and additional reasoning can also introduce a new arithmetic, scale, or evidence-selection error into an otherwise straightforward question.

This paper asks a narrower question than whether agents can answer financial questions: \emph{when is iterative agency worth its cost?} We study this question under controlled evidence access. Each system receives the same benchmark-provided, leakage-audited evidence corpus and the same permitted tools. This design does not measure end-to-end retrieval. Instead, it makes the orchestration decision identifiable: differences arise from how a system selects, organizes, and reasons over available evidence rather than from access to different documents.

Our method, \emph{Selective Finance Nexus}, begins with deterministic query analysis and evidence selection. A frozen gate then routes each example, before any answer generation, to either a one-call structured adjudicator or a bounded iterative ReAct process. The gate uses only observable, target-independent features of the question and evidence, such as length, numerical and temporal cues, lexical coverage, retrieval-score margin, and evidence dispersion. It is trained only on a development partition to predict the asymmetric event that ReAct succeeds while the static workflow fails. This formulation targets incremental value rather than generic difficulty.

We organize the study around four research questions:

\paragraph{RQ1.} Does deterministic-first orchestration improve the quality--cost frontier over direct prompting, structured one-call reasoning, and uniformly applied ReAct?
\paragraph{RQ2.} Which observable financial-question characteristics distinguish cases helped by iterative agency from cases solved more efficiently by a static workflow?
\paragraph{RQ3.} Can a development-frozen gate preserve the stronger system's accuracy while avoiding unnecessary iterative calls?
\paragraph{RQ4.} Are the conclusions stable across datasets, model families, and dataset-native evaluation protocols?

This development study presents a selective orchestration design, a leakage-audited paired protocol, and a measured negative gate outcome. The gate is scientifically consequential: the static and selective workflows did not satisfy the locked accuracy--cost criterion, so the disjoint 900-example final manifest was left unopened. We therefore preserve the confirmatory protocol as a later-study scaffold but make no held-out, cross-model, or end-to-end RAG claim here. The current artifact is suitable only for a workshop or protocol-focused negative-results report unless a redesigned method passes a new, independent development gate.
